\chapter{Conclusion}\label{chapter:conclusion}
Testing network protocol implementations for conformance with their specifications remains challenging due to protocol statefulness, complex dependencies across packet sequences, and large input spaces that are difficult to explore exhaustively.
Even minor deviations from specifications can lead to interoperability issues or security vulnerabilities, yet such deviations often escape detection by conventional testing techniques that rely on concrete test cases or incomplete oracles.

This thesis investigated how symbolic execution can be adapted to systematically and practically address the challenges of testing stateful protocol implementations.
It introduced a requirement-driven approach that uses symbolic execution to explore behaviors relevant to specification requirements, enabling the detection of violations that depend on specific corner-case values of protocol fields.
Building on this foundation, the thesis presented a monitor-based testing technique that externalizes requirement checking from the system under test, allowing protocol requirements to be expressed independently of specific implementations and reused across implementations and versions.
Finally, the thesis explored differential symbolic execution as a complementary approach that removes the need for a fully specified oracle by comparing observable behaviors across implementations, thereby enabling the detection of discrepancies that may indicate nonconformance issues even when explicit requirements have not been encoded.
As a comparative baseline, the thesis also empirically studied the use of coverage-guided fuzzing for testing protocol requirements and contrasted its effectiveness with symbolic execution in the context of CoAP implementations, highlighting the complementary strengths and limitations of the two techniques.

The proposed techniques were evaluated on multiple real-world implementations of the DTLS, CoAP, and QUIC protocols.
Across these studies, the techniques uncovered \emph{72} previously unknown specification violations.
These findings were validated in practice, with the majority of the reported issues fixed by developers.
The discovery of previously unknown bugs in widely deployed protocol implementations indicates that the proposed techniques are applicable beyond controlled academic settings and can be effectively used to analyze production-quality software.
Overall, the results demonstrate that symbolic execution, when guided by protocol requirements and combined with monitor-based and differential analysis, can systematically expose semantic deviations from protocol specifications that are difficult to detect using conventional testing techniques.